\documentclass[10pt,a4paper]{amsart}
\usepackage[margin=19mm]{geometry}
\usepackage{amsmath,amssymb,mathtools}
\usepackage[scr=rsfs,cal=euler]{mathalfa}
\usepackage{xfrac}
\usepackage[unicode,hidelinks,bookmarks=false]{hyperref}
\newcommand{\ie}{i.e.\ }
\newcommand{\cf}{cf.\ }
\newcommand{\resp}{respectively\ }
\newcommand{\Subsection}{\subsection}
\providecommand{\nobreakdash}{\nobreak\hbox{-}}
\setlength{\emergencystretch}{2em}
\multlinegap0pt
\usepackage{amssymb}
\usepackage{mathtools}
\usepackage{xfrac}
\usepackage{algorithm}\usepackage{algpseudocode}
\usepackage{xcolor}
\usepackage{float}
\newtheorem{lemma}{Lemma}
\title{Wedderburn decomposition of Twisted skew abelian group algebra}
\author{Alvaro Otero Sanchez}
\date{Source: arXiv:2609.07644v1 (CC BY 4.0)}
\begin{document}
\maketitle

\noindent\emph{Source and licence.} Wedderburn decomposition of Twisted skew abelian group algebra. Version arXiv:2609.07644v1 (CC BY 4.0), distributed by arXiv under the Creative Commons Attribution 4.0 International licence, http://creativecommons.org/licenses/by/4.0/, as recorded in the arXiv licence metadata for this exact version and shown on the abstract page https://arxiv.org/abs/2609.07644v1. Full licence notice: \texttt{licenses/arxiv-2609.07644-CC-BY-4.0.txt}.\par\smallskip

\noindent\textit{Editorial scope.} Counted unit: the article's lemma JinetaExtender (source0.tex lines 353--364). The article's complete original proof of it is delivered with it (source0.tex lines 365--411, one \texttt{proof} environment of 47 lines). No mathematics is written, repaired or re-derived: the statement and the proof are the article's own lines, in the article's own notation, and no retained line is substituted or re-flowed.  No further source-local context is needed: every cross-reference of every retained line resolves inside the two delivered spans.  Nothing was omitted from any retained span, so the package carries no omission record.  No retained line contains a citation, so the wrapper carries no bibliography and no bibliography entry.  No retained line contains a figure environment, an image-inclusion command or drawing code, so no image is delivered and the package declares no asset dependency.  Editorial formatting only, in addition: an AMS \texttt{amsart} preamble that restates the article's own declarations and macros used by the retained lines, namely the environments \texttt{lemma} on separate counters, together with the standard packages those lines need.  The retained environments print in this document as Lemma 1 in order of appearance, whereas the article prints the same items with the numbering of its own full document.  The complete native source of the article as distributed in the arXiv e-print for this exact version is bundled unchanged as \texttt{sources/209645/source0.tex}.
\begin{lemma} \label{JinetaExtender}
    Let $H=Q_{\hat{H}} Jin_\psi(H) \leq \hat{G}$, and let $C$ be a set of generators of the cyclic components of $H$ . Then $\psi:H \longrightarrow \mathbb{F}_q^*$ is determined up to $Q_{\hat{H}}$ by the image of the elements of $C$. In addition, a map $\psi: C \longrightarrow \mathbb{F}_q^*$ with $(\psi(a),a)\in H   \quad \forall a\in C$ can be extended to a Jineta of $H$ if and only if for all $a,b\in C$, we have that
    \begin{align} 
        \frac{\psi(a)}{\psi(a)^b} \cdot \frac{\psi(b)^a}{\psi(b)} \cdot \frac{\alpha(a,b)}{\alpha(b,a)}& \in Q_{\hat{H}} \\
       \prod_{k=1}^{ord(a)-1}\alpha(a,a^k) \prod_{k=0}^{ord(a)-1} \psi(a)^{a^k}& \in Q_{\hat{H}} 
    \end{align}
    or equivalently
        \begin{align} \label{EqDef}
        \left(\frac{\psi(a)}{\psi(a)^b} \cdot \frac{\psi(b)^a}{\psi(b)} \cdot \frac{\alpha(a,b)}{\alpha(b,a)} \right) ^ { |Q_{\hat{H}}|} =& 1 \\
       \left(\prod_{k=1}^{ord(a)-1}\alpha(a,a^k) \prod_{k=0}^{ord(a)-1} \psi(a)^{a^k}\right) ^ { |Q_{\hat{H}}|} =&  1 
    \end{align}
\end{lemma}

\begin{proof}
    First of all, we have that $(\psi(a),a), (\psi(b),b) \in H$ so
    \begin{equation}
        (\psi(a),a), (\psi(b),b) = ( \psi(a) \psi(b)^a, ab) \in H
    \end{equation}
    As we have that $(\psi(ab),ab) \in H$, then by the definition of $Q_{\hat{H}}$, we have that $ \frac{\psi(ab)}{\psi(a) \psi(b)^a \alpha(a,b)} \in Q_{\hat{H}}$. The first results come from induction on the number of products. 

    For the second part of the proposition, we know that if $\psi$ gives a Jineta of $H$, then it must satisfy the previous identity. In addition, as $G$ is abelian, we have that $\psi(ab) = \psi(ba)$. Therefore, we have that 
    \begin{equation}
        \psi(a) \psi(b)^a \alpha(a,b) Q_{\hat{H}} = \psi(b) \psi(a)^b \alpha(b,a)Q_{\hat{H}}
    \end{equation}
 Which leads to 
    \begin{equation}
         \frac{\psi(a)}{\psi(a)^b} \cdot \frac{\psi(b)^a}{\psi(b)} \cdot \frac{\alpha(a,b)}{\alpha(b,a)}  \in Q_{\hat{H}}
    \end{equation}
    As $\mathbb{F}_q^*$ is a cyclic group, $Q_{\hat{H}}$ is the only subgroup of order $ |Q_{\hat{H}}|$, and it has all elements $x\in \mathbb{F}_q^*$ such that $x^{|Q_{\hat{H}}|}$. For the second identity, we must have that $\pi_2((\psi(a),a)^{ord(a)}) = a^{ord(a)} = 1$, so $(\psi(a),a)^{ord(a)}) \in Q_{\hat{H}}$. Therefore, we have that 
    \begin{equation}
        (\psi(a),a)^{ord(a)} = \prod_{k=0}^{ord(a)-1}(\psi(a),a) = \left(\prod_{k=1}^{ord(a)-1}\alpha(a,a^k) \prod_{k=0}^{ord(a)-1} \psi(a)^{a^k}, 1\right) \in Q_{\hat{H}} 
    \end{equation}
    Reasoning as in the previous case, we have that 
    \begin{equation}
        \left(\prod_{k=1}^{ord(a)-1}\alpha(a,a^k) \prod_{k=0}^{ord(a)-1} \psi(a)^{a^k}\right) ^ { |Q_{\hat{H}}|} =1
    \end{equation}

    Now, for the if part of the second statement, suppose that those two equations are satisfied. Then, all elements in $H$ can be written as a word in terms of $C$, $ a=\prod_{i=1}^n a_i$ and we can define
    
    \begin{align*}
    \psi^* :  H &  \longrightarrow \mathbb{F}_q^* / Q_{\hat{H}} \\
        a=\prod_{i=1}^n a_i &\longmapsto \psi^*\left( \prod_{i=1}^h a_i \right) = \prod_{i=1}^{n} \left( \psi(a_i)^{\prod_{j=1}^{i-1} a_j} \alpha\left(\prod_{j=1}^{i-1} a_j, a_i\right)\right) \quad \mod Q_{\hat{H}} 
    \end{align*}
    
    To check that this map is well-defined, it suffices to show that $\psi^*(AabB) = \psi^*(AbaB)$ for all $A, B \in G$ and $a, b \in C$. The first identity of equation (\ref{EqDef}) provide that the word of length $2$ are well defined. By induction, lets suppose that the words up to $n-1$ are well defined. Then, we have that
    
    \begin{align*}
        \psi^*(AabB) & \equiv_{Q_{\hat{H}}} \psi(A)\psi(a)^A \psi(b)^{Aa} \psi(B)^{Aa} \alpha(A,a)\alpha(Aa,b)\alpha(Aab,B) \\ 
        & \equiv_{Q_{\hat{H}}} \psi(A) \left(\psi(a) \psi(b)^{a}\right)^{A} \psi(B)^{Aab} \alpha(A,a)\alpha(Aa,b)\alpha(Aab,B) \\ 
        & \equiv_{Q_{\hat{H}}} \psi(A) \left(\psi(b) \psi(a)^{b} \frac{\alpha(b,a)}{\alpha(a,b)}\right)^{A} \psi(B)^{Aba} \alpha(A,a)\alpha(Aa,b)\alpha(Aab,B) \\ 
        & \equiv_{Q_{\hat{H}}} \psi(A) \psi(b)^A \psi(a)^{Ab}  \psi(B)^{Aba} \alpha(A,a)\alpha(Aa,b)\alpha(Aab,B) \left(\frac{\alpha(b,a)}{\alpha(a,b)}\right)^{A} \\ 
        & \equiv_{Q_{\hat{H}}} \psi(A) \psi(b)^A \psi(a)^{Ab}  \psi(B)^{Aba} \alpha(a,b)^A\alpha(A,ab)\alpha(Aab,B) \left(\frac{\alpha(b,a)}{\alpha(a,b)}\right)^{A}\\
        & \equiv_{Q_{\hat{H}}} \psi(A) \psi(b)^A \psi(a)^{Ab}  \psi(B)^{Aba} \alpha(a,b)^A\alpha(A,ba)\alpha(Aab,B)\left(\frac{\alpha(b,a)}{\alpha(a,b)}\right)^{A} \\ 
        & \equiv_{Q_{\hat{H}}} \psi(A) \psi(b)^A \psi(a)^{Ab}  \psi(B)^{Aba} \frac{\alpha(a,b)^A}{\alpha(b,a)^A}\alpha(A,b) \alpha(Ab,a)\alpha(Aab,B) \left(\frac{\alpha(b,a)}{\alpha(a,b)}\right)^{A}\\ 
        & \equiv_{Q_{\hat{H}}} \psi(A) \psi(b)^A \psi(a)^{Ab}  \psi(B)^{Aba} \alpha(A,b)\alpha(Ab,a)\alpha(Aba,B) \left(\frac{\alpha(a,b)}{\alpha(b,a)}\right)^{A} \\ 
        & \equiv_{Q_{\hat{H}}} \psi(AbaB)
    \end{align*}

    Therefore, we can define $\psi(a)$ as a representative of $\psi^*(a)$. Finally, it is trivial to see that $\{(\psi(a),a) ; a \in \pi_2(H)\} \subset H$ and therefore it define a Jineta for $H$
\end{proof}

\end{document}
